% =============================================================================
% Section: Experiments
% =============================================================================
% Required packages (usually already included in the CVPR template):
%   \usepackage{booktabs}
%   \usepackage{multirow}
%   \usepackage{graphicx}
%   \usepackage{amsmath}
%   \usepackage{float}       % only if [H] is used
%
% =============================================================================

\section{Experiments}
\label{sec:experiments}

We evaluate \textbf{RepairOPD} from four complementary perspectives:
(1) whether mask-policy OPD enables effective interactive mask refinement;
(2) whether repair-oriented training preserves one-pass conversational segmentation;
(3) whether the proposed Data-Policy and Mask-Policy components are individually effective and complementary; and
(4) how the learned repair behavior emerges under different error states and affects inference efficiency.
We first describe the experimental setup, followed by the main results, component ablations, and detailed analysis.


% -----------------------------------------------------------------------------
\subsection{Experimental Setup}
\label{sec:exp_setup}

\paragraph{Benchmarks.}
We evaluate RepairOPD on three groups of benchmarks.

\noindent\textbf{ConVerSeg repair protocol.}
Our primary evaluation focuses on conversational mask repair.
Each sample consists of an image $I$, a conversational prompt $p$, and a ground-truth mask $M$.
We hold out a fixed subset from the conversational Stage-2 data mixture for all repair experiments.

To evaluate different types of flawed predictions, we consider two complementary bad-mask protocols.
First, in the \textbf{on-policy} setting, the initial mask is generated by the student's own text-only prediction,
i.e., the model first predicts $q_s^{1}=\pi_{\theta}(I,p,\emptyset)$ and then receives this prediction as the mask state for refinement.
This setting directly approximates interactive deployment, where users revise the model's own initial prediction.

Second, in the \textbf{corrupted-GT} setting, the initial mask is obtained by applying controlled stochastic corruptions to the ground-truth mask.
We consider perturbations such as erosion, dilation, missing components, false regions, boundary jitter, and spatial displacement.
This protocol provides controlled error states and enables a direct comparison with the Path-1 correction teacher.
Both protocols use identical image--prompt--target-mask pairs; only the source of the flawed mask differs.

\noindent\textbf{ConVerSeg benchmark.}
We further evaluate one-pass conversational image segmentation on the official ConVerSeg evaluation benchmark.
Following the original protocol, models receive only the image and conversational text without a mask prompt at inference time.
We report results on both the \textbf{SAM-seeded} and \textbf{human-annotated} splits across five concept families:
entities, spatial relations, affordances, relations \& events, and physics \& safety.
This evaluation measures whether repair-oriented training preserves the model's general conversational grounding capability.

\noindent\textbf{Referring and reasoning segmentation.}
Following prior work, we optionally evaluate zero-shot transfer on RefCOCO, RefCOCO+, RefCOCOg, and ReasonSeg.
These benchmarks provide additional evidence on whether the learned repair behavior affects standard referring-expression grounding and more complex reasoning-based segmentation.
No task-specific fine-tuning is performed on these datasets.


\paragraph{Baselines.}
We compare RepairOPD with the following baselines and references.

\noindent\textbf{Stage-2 ConVerSeg student (Ours-init).}
Our primary baseline is the Stage-2 conversational segmentation checkpoint used to initialize RepairOPD.
It is trained with the same conversational data mixture but without mask-policy OPD and without repair-oriented distillation.
At repair time, the same mask-conditioned pathway is exposed to the baseline, allowing us to determine whether standard conversational segmentation training alone is sufficient for effective mask correction.

\noindent\textbf{RepairOPD (Ours).}
RepairOPD is our full model, initialized from the Stage-2 checkpoint and trained with the proposed Mask-Policy OPD and Data-Policy OPD.
The Path-1 correction teacher is frozen during student training.
All trainable student components are jointly optimized.

\noindent\textbf{Path-1 correction teacher (Oracle).}
The correction teacher is trained explicitly to recover the ground-truth mask from corrupted mask states.
It provides the teacher supervision for Mask-Policy OPD and serves as an upper reference for repair quality.
The teacher is not used during inference.

\noindent\textbf{ConVerSeg-Net (Reference).}
We additionally report published results from ConVerSeg-Net~\cite{sahoo2026convseg}, and use the released checkpoint as an external reference where evaluation under the corresponding protocol is available.
Because our Stage-2 initialization uses a reduced training budget relative to the complete ConVerSeg curriculum, this reference is not treated as a strictly matched training baseline.

\noindent\textbf{External segmentation methods.}
For broader comparison on one-pass conversational segmentation, we optionally include representative MLLM- and SAM-based methods reported in the ConVerSeg study, such as LISA~\cite{lai2024lisa}, UniLSeg~\cite{unilseg}, EVF-SAM~\cite{evfsam}, and Seg-Zero~\cite{segzero}.
These methods are included only as external references and are not retrained for mask repair.


\paragraph{Evaluation Metrics.}
We use task-specific metrics for one-pass segmentation and interactive mask repair.

For \textbf{one-pass conversational segmentation}, we report generalized intersection-over-union (gIoU) as the primary metric following prior work.
Cumulative IoU (cIoU) is additionally reported in the Appendix.

For \textbf{mask repair}, we report four complementary metrics.
Let $M_b$ denote the initial flawed mask, $M_r$ the repaired mask, and $M$ the ground-truth mask.

\begin{itemize}
    \item \textbf{Text IoU} measures the IoU between the student's text-only prediction and the ground-truth mask:
    \[
    \mathrm{IoU}_{\mathrm{text}}
    =
    \mathrm{IoU}(q_s^{1},M).
    \]

    \item \textbf{IoU$_{\mathrm{after}}$} measures the absolute quality of the repaired prediction:
    \[
    \mathrm{IoU}_{\mathrm{after}}
    =
    \mathrm{IoU}(M_r,M).
    \]
    We use this metric as the primary measure of post-repair segmentation quality.

    \item \textbf{Recovery Gain} measures the improvement over the flawed initial hypothesis:
    \[
    \mathrm{Gain}
    =
    \mathrm{IoU}_{\mathrm{after}}
    -
    \mathrm{IoU}_{\mathrm{before}},
    \]
    where
    $\mathrm{IoU}_{\mathrm{before}}
    =
    \mathrm{IoU}(M_b,M)$.
    A positive gain indicates that refinement improves the initial mask on average.

    \item \textbf{Fraction with Gain $>0$} measures the percentage of samples for which repair strictly improves over the initial mask:
    \[
    \mathrm{Frac.}_{>0}
    =
    \frac{1}{N}
    \sum_{i=1}^{N}
    \mathbb{I}
    \left[
    \mathrm{Gain}_i>0
    \right].
    \]
    This metric measures refinement stability rather than only average improvement.
\end{itemize}

Unless otherwise specified, predicted masks are binarized using a threshold of $0.5$.
Ground-truth erosion follows the corresponding ConVerSeg evaluation protocol where applicable.


\paragraph{Implementation Details.}
The student model uses \textbf{SAM2-Large}~\cite{ravi2024sam2} as the image encoder and mask decoder backbone and \textbf{Qwen2.5-VL-3B}~\cite{qwen25vl} as the prompt encoder.
The Qwen2.5-VL parameters are adapted using LoRA~\cite{hu2022lora} with rank $16$ and scaling factor $\alpha=32$.

RepairOPD is initialized from the Stage-2 conversational checkpoint.
The Path-1 correction teacher uses Qwen2.5-VL-7B with SAM2 and is trained using corrupted ground-truth states.
The teacher is frozen before student OPD training.

For the student, the SAM2 mask decoder, MaskStateEncoder, and language-side LoRA adapters are jointly optimized.
Mask fusion uses $\gamma=0.25$.
The KD temperature is set to $T\approx1.5$, and the default probability of sampling on-policy states is $\rho\approx0.5$, with the remaining samples drawn from the corrupted-GT protocol.

The student is optimized using AdamW~\cite{loshchilov2019adamw} with mixed-precision training, batch size $1$ and gradient accumulation.
The training data follows the conversational Stage-2 mixture, augmented by hard samples produced by Data-Policy OPD.
No additional human repair annotations are introduced.

At inference time, the model supports both text-only grounding and mask-conditioned repair.
The two modes share cached SAM image embeddings, allowing the repair pathway to reuse the image representation computed during the initial prediction.
Complete optimization hyperparameters, corruption configurations, training schedule, and hardware settings are provided in the Appendix.


% -----------------------------------------------------------------------------
\subsection{Main Results}
\label{sec:main_results}

\paragraph{Mask Repair.}
Table~\ref{tab:repair_main} summarizes repair performance under both on-policy and corrupted-GT protocols.

The comparison with the Stage-2 initialization directly measures whether repair-oriented training enables useful correction beyond standard conversational segmentation.
The on-policy protocol evaluates the realistic deployment scenario in which the model must refine its own initial prediction.
In contrast, the corrupted-GT protocol provides controlled and potentially more severe error states, allowing us to examine the model's intrinsic correction capability independently of its initial grounding errors.

We evaluate both absolute post-repair quality and improvement relative to the initial flawed mask.
In particular, $\mathrm{IoU}_{\mathrm{after}}$ measures the final segmentation quality, while Recovery Gain measures the ability to recover from the specific error state.
The fraction of samples with positive gain further indicates whether the learned refinement behavior is consistently beneficial rather than being dominated by a small number of large improvements.

The Path-1 teacher provides an upper reference for the repair task because it is explicitly trained on corrupted mask states.
The gap between the teacher and RepairOPD therefore measures the extent to which correction behavior can be transferred to the deployable student without requiring the teacher at inference time.


\begin{table*}[t]
    \centering
    \caption{
    \textbf{Mask repair results} on the held-out repair subset.
    Results are reported under both on-policy and corrupted-GT protocols.
    $\mathrm{Gain}$ denotes $\mathrm{IoU}_{\mathrm{after}}-\mathrm{IoU}_{\mathrm{before}}$.
    \textbf{Bold} indicates the best student performance.
    The Path-1 teacher is used only during training and serves as an oracle reference.
    }
    \label{tab:repair_main}
    \small
    \setlength{\tabcolsep}{5pt}
    \begin{tabular}{l cc cc cc cc}
        \toprule
        & \multicolumn{4}{c}{\textbf{On-policy}} &
        \multicolumn{4}{c}{\textbf{Corrupted GT}} \\
        \cmidrule(lr){2-5}
        \cmidrule(lr){6-9}
        \textbf{Method}
        & Text IoU
        & IoU$_{\mathrm{after}}$
        & Gain
        & Frac.$\!>\!0$
        & IoU$_{\mathrm{before}}$
        & IoU$_{\mathrm{after}}$
        & Gain
        & Frac.$\!>\!0$ \\
        \midrule
        Stage-2 (Ours-init)
        & --
        & --
        & --
        & --
        & --
        & --
        & --
        & -- \\

        RepairOPD (Ours)
        & \textbf{--}
        & \textbf{--}
        & \textbf{--}
        & \textbf{--}
        & --
        & \textbf{--}
        & \textbf{--}
        & \textbf{--} \\

        Path-1 Teacher (Oracle)
        & --
        & --
        & --
        & --
        & --
        & --
        & --
        & -- \\
        \bottomrule
    \end{tabular}
\end{table*}


\paragraph{One-Pass Conversational Segmentation.}
Table~\ref{tab:convseg_main} evaluates whether repair-oriented training preserves standard conversational image segmentation.

All models are evaluated under the text-only protocol without mask prompts.
We report gIoU on both the SAM-seeded and human-annotated splits of ConVerSeg.
The Stage-2 initialization provides the direct reference for measuring the effect of RepairOPD training, while ConVerSeg-Net serves as an external reference trained under the full original curriculum.

Our primary objective is not to replace the original one-pass ConVerSeg training recipe, but to determine whether introducing repair-oriented supervision damages the model's existing grounding capability.
A comparable or improved one-pass result together with improved repair performance would indicate that mask-conditioned correction can be learned without sacrificing standard conversational segmentation.


\begin{table}[t]
    \centering
    \caption{
    \textbf{ConVerSeg gIoU} under the standard text-only protocol.
    ``All'' aggregates the SAM-seeded and human-annotated splits.
    ConVerSeg-Net is included as an external reference.
    }
    \label{tab:convseg_main}
    \small
    \setlength{\tabcolsep}{5pt}
    \begin{tabular}{lccc}
        \toprule
        \textbf{Method}
        & SAM-seeded
        & Human-ann.
        & All \\
        \midrule
        Stage-2 (Ours-init)
        & --
        & --
        & -- \\

        RepairOPD (Ours)
        & \textbf{--}
        & \textbf{--}
        & \textbf{--} \\

        \midrule

        ConVerSeg-Net-3B (Ref.)~\cite{sahoo2026convseg}
        & --
        & --
        & -- \\
        \bottomrule
    \end{tabular}
\end{table}


\begin{table}[t]
    \centering
    \caption{
    \textbf{Per-concept gIoU} on the ConVerSeg human-annotated split.
    Evaluation follows the text-only protocol.
    }
    \label{tab:convseg_concept}
    \small
    \setlength{\tabcolsep}{3.5pt}
    \begin{tabular}{lccccc}
        \toprule
        \textbf{Method}
        & Ent.
        & Spat.
        & Aff.
        & Rel./Ev.
        & Phys./Saf. \\
        \midrule
        Stage-2 (Ours-init)
        & --
        & --
        & --
        & --
        & -- \\

        RepairOPD (Ours)
        & \textbf{--}
        & \textbf{--}
        & \textbf{--}
        & \textbf{--}
        & \textbf{--} \\
        \bottomrule
    \end{tabular}
\end{table}


\paragraph{Transferability.}
Table~\ref{tab:refcoco} reports zero-shot transfer results on RefCOCO, RefCOCO+, RefCOCOg, and ReasonSeg.
No task-specific fine-tuning is performed.

These experiments are not intended to establish state-of-the-art performance on the individual benchmarks.
Instead, they examine whether the repair-oriented training procedure changes the general segmentation behavior of the Stage-2 model.
Comparable performance between the initialization and RepairOPD indicates that the proposed repair pathway does not substantially compromise general referring-expression or reasoning-based segmentation.


\begin{table*}[t]
    \centering
    \caption{
    \textbf{Zero-shot transfer} on RefCOCO/+/g and ReasonSeg.
    Models are not fine-tuned on these benchmarks.
    }
    \label{tab:refcoco}
    \small
    \setlength{\tabcolsep}{4pt}
    \begin{tabular}{lcccccccc}
        \toprule
        \textbf{Method}
        & \multicolumn{3}{c}{RefCOCO}
        & \multicolumn{3}{c}{RefCOCO+}
        & RefCOCOg
        & ReasonSeg \\
        \cmidrule(lr){2-4}
        \cmidrule(lr){5-7}
        \cmidrule(lr){8-8}
        \cmidrule(lr){9-9}
        & val & testA & testB
        & val & testA & testB
        & val & test \\
        \midrule
        Stage-2 (Ours-init)
        & -- & -- & --
        & -- & -- & --
        & -- & -- \\

        RepairOPD (Ours)
        & \textbf{--} & \textbf{--} & \textbf{--}
        & \textbf{--} & \textbf{--} & \textbf{--}
        & \textbf{--} & \textbf{--} \\

        ConVerSeg-Net-3B (Ref.)~\cite{sahoo2026convseg}
        & -- & -- & --
        & -- & -- & --
        & -- & -- \\
        \bottomrule
    \end{tabular}
\end{table*}


\paragraph{Qualitative Results.}
Figure~\ref{fig:qualitative} presents representative conversational mask repair examples.
We compare the initial text-only prediction, the flawed mask state, the RepairOPD refinement, and the ground-truth mask.

The examples cover several characteristic error patterns, including over-segmentation, under-segmentation, spatial displacement, missing regions, and distractor inclusion.
RepairOPD is expected to use the flawed mask as an explicit spatial state while relying on the conversational prompt to determine which regions should be preserved, removed, or recovered.

We additionally visualize failure cases.
These cases are particularly informative when the initial mask is severely corrupted or when the prompt requires abstract functional or physical reasoning.
Additional comparisons with the Path-1 teacher are provided in the Appendix.


\begin{figure}[t]
    \centering
    % \includegraphics[width=\linewidth]{figures/repair_qualitative.pdf}
    \fbox{
        \parbox[c][1.75in][c]{0.95\linewidth}{
            \centering
            \small
            Placeholder: qualitative repair examples\\
            Initial mask / flawed mask / RepairOPD / GT\\
            Success cases + failure cases
        }
    }
    \caption{
    \textbf{Qualitative mask repair examples.}
    RepairOPD receives the same conversational prompt and flawed mask as input and refines the prediction toward the intended target region.
    }
    \label{fig:qualitative}
\end{figure}


% -----------------------------------------------------------------------------
\subsection{Ablation Studies}
\label{sec:ablations}

\paragraph{Effect of Mask-Policy OPD.}
We first isolate the contribution of Mask-Policy OPD by fixing the training data distribution.
The key comparison is between standard mask-conditioned training and the proposed student-induced state formulation.

Specifically, we compare:
(1) the Stage-2 initialization without repair training;
(2) training with fixed or corrupted mask states;
(3) training with historical student states;
(4) training with randomly corrupted states; and
(5) the proposed current-state on-policy Mask-Policy OPD.

This comparison tests whether the correction state should be generated from the student's current prediction rather than from an externally predefined corruption process.
The experiment directly evaluates the state-matching hypothesis underlying Mask-Policy OPD.


\begin{table}[t]
    \centering
    \caption{
    \textbf{Mask-Policy OPD ablation.}
    The data distribution is fixed while different mask-state construction strategies are compared.
    }
    \label{tab:ablation_mask_policy}
    \small
    \setlength{\tabcolsep}{4pt}
    \begin{tabular}{lccc}
        \toprule
        \textbf{Variant}
        & State Source
        & IoU$_{\mathrm{after}}$
        & Gain \\
        \midrule
        Stage-2
        & None
        & --
        & -- \\

        Off-policy KD
        & GT / Fixed
        & --
        & -- \\

        Historical-State KD
        & Historical
        & --
        & -- \\

        Random-State KD
        & Random corrupt.
        & --
        & -- \\

        Mask-Policy OPD
        & \textbf{Current}
        & \textbf{--}
        & \textbf{--} \\
        \bottomrule
    \end{tabular}
\end{table}


\paragraph{Effect of Data-Policy OPD.}
We next isolate the contribution of the outer Data-Policy OPD while keeping the student architecture and mask-state mechanism fixed.

We compare a fixed data distribution, uniform random augmentation, and the proposed error-driven data policy under the same additional-sample budget.
The error-driven policy uses the student error profile to identify underperforming concept families and then generates targeted hard prompts through semantic-preserving rewrites, minimal semantic pairs, compositional constraints, functional or physical counterfactuals, and negative prompts.

The comparison evaluates whether dynamically allocating new samples according to the student's current weaknesses provides greater training value than increasing the dataset uniformly.


\begin{table}[t]
    \centering
    \caption{
    \textbf{Data-Policy OPD ablation.}
    All variants use the same additional-sample budget.
    }
    \label{tab:ablation_data_policy}
    \small
    \setlength{\tabcolsep}{4pt}
    \begin{tabular}{lccc}
        \toprule
        \textbf{Variant}
        & Data Policy
        & gIoU
        & Repair Gain \\
        \midrule
        Fixed Data
        & Fixed
        & --
        & -- \\

        Random-Augment
        & Random
        & --
        & -- \\

        Data-Policy OPD
        & \textbf{Error-driven}
        & \textbf{--}
        & \textbf{--} \\
        \bottomrule
    \end{tabular}
\end{table}


\paragraph{Synergy of Dual-Level OPD.}
The previous experiments isolate the two levels independently.
We now evaluate whether their combination provides complementary benefits.

Table~\ref{tab:ablation_dual} compares:
(1) the Stage-2 baseline,
(2) Data-Policy OPD only,
(3) Mask-Policy OPD only, and
(4) the full DL-OPD framework.

Data-Policy OPD operates at the distribution level and focuses additional training capacity on persistent semantic, relational, functional, physical, negative-prompt, and geometric weaknesses.
Mask-Policy OPD operates at the state level and explicitly trains the student to correct its own flawed predictions.
The full framework combines these two mechanisms.


\begin{table}[t]
    \centering
    \caption{
    \textbf{Synergy of dual-level OPD.}
    Data-Policy and Mask-Policy operate at the data-distribution and mask-state levels, respectively.
    }
    \label{tab:ablation_dual}
    \small
    \setlength{\tabcolsep}{4pt}
    \begin{tabular}{lcccc}
        \toprule
        \textbf{Variant}
        & Data Policy
        & Mask State
        & gIoU
        & Repair Gain \\
        \midrule
        Stage-2
        & Fixed
        & None
        & --
        & -- \\

        Data-Policy only
        & Error-driven
        & None
        & --
        & -- \\

        Mask-Policy only
        & Fixed
        & Current
        & --
        & -- \\

        \textbf{DL-OPD}
        & \textbf{Error-driven}
        & \textbf{Current}
        & \textbf{--}
        & \textbf{--} \\
        \bottomrule
    \end{tabular}
\end{table}


\paragraph{Mask Fusion and Bad-Mask Mixture.}
We separately study two design choices that control how mask states are incorporated into the student.

First, we vary the mask fusion coefficient $\gamma$ in the dense prompt residual.
When $\gamma=0$, no mask information is injected into the dense prompt representation.
Increasing $\gamma$ strengthens the spatial conditioning provided by the flawed mask.
The experiment measures the trade-off between conversational semantic information and spatial mask conditioning.

Second, we vary the bad-mask mixture probability $\rho$ used during OPD.
A low $\rho$ exposes the student primarily to controlled corrupted-GT states, while a high $\rho$ emphasizes its own on-policy predictions.
This experiment determines how deployment-aligned errors and severe synthetic errors should be balanced during training.


\begin{table}[t]
    \centering
    \caption{
    \textbf{Mask fusion scale} $\gamma$ ablation.
    }
    \label{tab:ablation_gamma}
    \small
    \begin{tabular}{lccc}
        \toprule
        $\gamma$
        & IoU$_{\mathrm{after}}$
        & Gain
        & gIoU \\
        \midrule
        0.00
        & --
        & --
        & -- \\

        0.25
        & \textbf{--}
        & \textbf{--}
        & \textbf{--} \\

        0.50
        & --
        & --
        & -- \\
        \bottomrule
    \end{tabular}
\end{table}


\begin{table}[t]
    \centering
    \caption{
    \textbf{Bad-mask mixture} $\rho$ ablation.
    $\rho$ denotes the probability of sampling on-policy student states during training.
    }
    \label{tab:ablation_rho}
    \small
    \setlength{\tabcolsep}{4pt}
    \begin{tabular}{lcccc}
        \toprule
        $\rho$
        & \multicolumn{2}{c}{On-policy eval.}
        & \multicolumn{2}{c}{Corrupted-GT eval.} \\
        \cmidrule(lr){2-3}
        \cmidrule(lr){4-5}
        & IoU$_{\mathrm{after}}$
        & Gain
        & IoU$_{\mathrm{after}}$
        & Gain \\
        \midrule
        0.0
        & --
        & --
        & --
        & -- \\

        0.5
        & \textbf{--}
        & \textbf{--}
        & \textbf{--}
        & \textbf{--} \\

        1.0
        & --
        & --
        & --
        & -- \\
        \bottomrule
    \end{tabular}
\end{table}


\paragraph{Teacher Distillation.}
We investigate the contribution of teacher supervision by removing the teacher KD objective while retaining the remaining supervised losses.

The full model receives soft correction targets from the frozen Path-1 teacher.
Without teacher KD, the student is optimized only using the available hard supervision and auxiliary objectives.
Comparing the two variants determines whether teacher predictions provide correction information beyond the ground-truth target itself.

We further compare the student against the Path-1 teacher under the corrupted-GT protocol to quantify the remaining teacher--student gap.


\begin{table}[t]
    \centering
    \caption{
    \textbf{Teacher distillation ablation} on the repair subset.
    }
    \label{tab:ablation_teacher}
    \small
    \setlength{\tabcolsep}{4pt}
    \begin{tabular}{lcccc}
        \toprule
        \textbf{Variant}
        & IoU$_{\mathrm{after}}$
        & Gain
        & Frac.$\!>\!0$
        & gIoU \\
        \midrule
        w/o Teacher KD
        & --
        & --
        & --
        & -- \\

        RepairOPD
        & \textbf{--}
        & \textbf{--}
        & \textbf{--}
        & \textbf{--} \\
        \bottomrule
    \end{tabular}
\end{table}

